\documentclass[11pt]{article}

\usepackage[margin=1in]{geometry}
\usepackage[utf8]{inputenc}
\usepackage[T1]{fontenc}
\usepackage{lmodern}
\usepackage{xcolor}
\usepackage{titlesec}
\usepackage{enumitem}
\usepackage{parskip}
\usepackage[colorlinks=true, linkcolor=blue!50!black, citecolor=blue!50!black, urlcolor=blue!50!black]{hyperref}

\usepackage[backend=biber, style=authoryear, sorting=nyt]{biblatex}
\addbibresource{references.bib}

\titleformat{\section}{\Large\bfseries}{\thesection}{0.6em}{}
\titleformat{\subsection}{\large\bfseries}{\thesubsection}{0.6em}{}
\titlespacing*{\subsection}{0pt}{1.4em}{0.4em}

\newcommand{\charname}[1]{\textbf{#1}}

\title{\textbf{A Reader's Companion to \emph{1929}}\\[0.3em]
\large Andrew Ross Sorkin's \emph{1929: Inside the Greatest Crash in Wall Street History \\ and How It Shattered a Nation}}
\author{Prepared by Claude (Sonnet 5), Anthropic \\ \small for \href{mailto:icos.atropa@gmail.com}{a listener of the audiobook narration}}
\date{August 24, 2026}

\begin{document}

\maketitle

\begin{center}
\begin{minipage}{0.85\textwidth}
\small
\textit{Note on this document.} This companion is a listening aid, not a substitute for
the book. It was assembled entirely from publicly available secondary sources ---
publisher copy, book reviews, press interviews, and reader summaries --- rather than
from the text of \emph{1929} itself, and every claim below is cited to one of those
public sources in the References section. Because the goal is to help a listener keep
names straight, character and institution sketches describe who people \emph{were} and
how they connect to one another; they deliberately omit how each person's individual
story concludes. The chapter-by-chapter section is a conceptual reconstruction of the
book's chronological arc, not a transcription of its actual chapter titles: reviewers
note that Sorkin structures the book as a sequence of date-headed chapters (beginning
February 1, 1929) grouped into two parts, rather than as topically titled chapters, and
that precise chapter-by-chapter titles are not published online
\autocite{wikipedia1929book,supersummary1929}. This document is not authorized,
reviewed, or endorsed by the author or publisher.
\end{minipage}
\end{center}

\vspace{1em}

\section*{Overview}

\emph{1929} (Viking, 2025) is Andrew Ross Sorkin's narrative history of the Wall
Street crash of 1929, its long buildup, and its aftermath through the New Deal-era
banking reforms \autocite{sorkin2025book,prh1929,goodreads1929}. Sorkin, best known for his account
of the 2008 financial crisis in \emph{Too Big to Fail}, spent roughly eight years
researching the book, drawing on archival material including previously unpublished
Federal Reserve Bank of New York board records \autocite{axios1929,realclearmarkets1929}.
Rather than telling the crash as a story of charts and abstractions, the book is built
around a cast of interlocking real people --- bankers, speculators, regulators, and
politicians --- whose personal ambitions and rivalries shaped the boom, the crash, and
the decade of reform that followed \autocite{airmail1929,csmonitor1929}. Reviewers
describe the book as unfolding in two parts: the first tracing the speculative boom
from roughly February through October 1929, and the second following the crash's
political and regulatory aftermath \autocite{cato1929,dailyeconomy1929}.

\section{Cast of Characters}

The figures below are the major, named individuals the book follows closely across
its narrative, drawn from publisher and reviewer descriptions of the book's
principal cast \autocite{supersummary1929,booksthatslay1929,littlerbooks1929}. Wilson
and Coolidge, who appear only briefly as outgoing political context, are omitted at
the reader's request.

\subsection{Charles E. Mitchell}
Chairman of National City Bank (a predecessor of Citibank) and known as ``Sunshine
Charlie'' for his relentless public optimism, Mitchell pioneered the mass marketing
of stocks to ordinary Americans and used his bank's resources to keep credit flowing
to the market even when the Federal Reserve wanted to restrain it
\autocite{airmail1929,supersummary1929}. His running conflict with the Federal
Reserve's leadership is one of the book's central threads connecting Wall Street to
Washington.

\subsection{Thomas W. Lamont}
The senior partner at J.P.\ Morgan \& Co.\ and the most influential private banker of
the era, Lamont acted as an informal statesman for Wall Street, coordinating with
fellow bankers and government officials to manage the market's mood
\autocite{supersummary1929,dailyeconomy1929}. He is a connective figure between the
House of Morgan, the New York Stock Exchange, and the Hoover administration.

\subsection{J.P.\ Morgan Jr.}
Head of the House of Morgan and heir to his father's banking dynasty, Morgan
represents the old-guard establishment of American finance, whose firm's private
influence over markets and government rivaled that of the Federal Reserve itself
\autocite{supersummary1929}. He works closely with Lamont throughout the book.

\subsection{Richard Whitney}
A vice president (and later president) of the New York Stock Exchange, Whitney
became the public face of Wall Street's efforts to project confidence on the trading
floor during the market's most volatile days, acting on behalf of the leading
bankers' consortium \autocite{supersummary1929,csmonitor1929}. He works in close
coordination with Lamont and the Morgan-led banking group.

\subsection{Jesse Livermore}
A legendary independent stock trader famous for making and losing several fortunes,
Livermore operated outside the banking establishment, trading on his own instincts
about market psychology rather than on inside information or institutional loyalty
\autocite{airmail1929,supersummary1929}. He functions in the narrative as a
skeptical outsider watching the same boom that bankers like Mitchell were fueling.

\subsection{William C. Durant}
The founder of General Motors, Durant reinvented himself in the late 1920s as one of
Wall Street's most aggressive stock market operators, organizing pools of investors
to move prices in favored stocks \autocite{supersummary1929}. His path connects
General Motors, the speculative pool culture of the exchange, and fellow GM figure
John J.\ Raskob.

\subsection{John J.\ Raskob}
A General Motors and DuPont executive who also chaired the Democratic National
Committee, Raskob was a leading public evangelist for stock ownership, urging
ordinary Americans that ``everybody ought to be rich'' through the market
\autocite{supersummary1929,littlerbooks1929}. He links the corporate world of
General Motors to both Wall Street speculation and national politics.

\subsection{Michael J.\ Meehan}
A New York Stock Exchange specialist and stock promoter known for running trading
pools in popular speculative stocks such as Radio Corporation of America (RCA),
Meehan exemplifies the era's loosely regulated market-manipulation practices
\autocite{supersummary1929}. His activities intersect with the same pool culture
that drew in Durant and Raskob.

\subsection{Herbert Hoover}
Inaugurated president in March 1929, Hoover inherited a market boom he privately
distrusted and then had to respond to its collapse, caught between his instincts for
voluntary, business-led solutions and mounting pressure for federal action
\autocite{supersummary1929,csmonitor1929}. His administration's dealings with bankers
like Lamont and Mitchell are a recurring axis of the book.

\subsection{Andrew W.\ Mellon}
A wealthy financier and industrialist serving as Treasury Secretary, Mellon
represents the fusion of private fortune and public economic policy that
characterized the era, shaping the government's hands-off approach to the boom
\autocite{supersummary1929}. He works alongside Hoover through the crash's early
aftermath.

\subsection{George L.\ Harrison}
Governor of the Federal Reserve Bank of New York, Harrison inherited the internal
Federal Reserve debates over whether and how to curb speculative lending, and dealt
directly with bankers such as Mitchell whose actions the Fed was trying to restrain
\autocite{supersummary1929}. He is the book's throughline for the Federal Reserve's
institutional perspective.

\subsection{Franklin D.\ Roosevelt}
Governor of New York during the crash and elected president in 1932, Roosevelt
becomes central to the book's second half as the political figure whose
administration presided over the era's regulatory response to Wall Street's excesses
\autocite{supersummary1929,libraryjournal1929}. He connects to Senator Carter Glass
and investigator Ferdinand Pecora as reform-era figures.

\subsection{Carter Glass}
A U.S.\ Senator and former Treasury Secretary, Glass was a leading architect of
banking reform legislation in the crash's aftermath, working to redraw the boundary
between commercial and investment banking \autocite{libraryjournal1929}. His
legislative efforts connect Washington's political response to the practices of
bankers like Mitchell and Lamont profiled earlier in the book.

\subsection{Ferdinand Pecora}
Chief counsel to the U.S.\ Senate Committee on Banking and Currency, Pecora led the
high-profile investigative hearings into Wall Street practices that followed the
crash, questioning many of the book's banking figures directly about their conduct
during the boom \autocite{libraryjournal1929,dailyeconomy1929}. His hearings serve
as the narrative bridge between the crash and the reforms that followed it.

\section{Cast of Institutions}

\subsection{National City Bank}
A forerunner of Citibank and, under Charles Mitchell's leadership, one of the most
aggressive commercial banks in extending credit for stock speculation and marketing
securities directly to retail customers \autocite{airmail1929,supersummary1929}. Its
practices put it repeatedly at odds with the Federal Reserve.

\subsection{J.P.\ Morgan \& Co.}
The most powerful private bank on Wall Street, led by Thomas Lamont and J.P.\ Morgan
Jr., whose partners acted as informal coordinators of the banking community's
response to market stress \autocite{supersummary1929}. The firm's influence
connects private finance directly to the New York Stock Exchange and to Washington.

\subsection{New York Stock Exchange (NYSE)}
The physical and symbolic center of American securities trading, where figures like
Richard Whitney and Michael Meehan operated, and whose self-regulated trading floor
practices are a recurring subject of the book \autocite{supersummary1929}. It is the
stage on which the boom's speculative excesses and the crash itself play out.

\subsection{The Federal Reserve System}
Comprising the Board in Washington and regional banks such as the Federal Reserve
Bank of New York (led by George Harrison), the Fed grappled throughout the boom with
whether and how to restrain speculative lending by banks like National City
\autocite{supersummary1929}. Its internal debates connect monetary policy to the
conduct of individual bankers profiled elsewhere in the book.

\subsection{General Motors}
The automaker at the center of two major characters' fortunes --- founder William C.\
Durant and executive John J.\ Raskob --- General Motors' stock was itself one of the
era's most actively speculated-upon issues, tying industrial America directly to the
market's mania \autocite{supersummary1929}.

\subsection{The White House and U.S.\ Treasury}
The executive branch under Herbert Hoover and Treasury Secretary Andrew Mellon, and
later under Franklin Roosevelt, whose evolving posture toward Wall Street --- from
voluntary cooperation to active regulation --- structures the political half of the
book's narrative \autocite{supersummary1929,csmonitor1929}.

\subsection{U.S.\ Senate Committee on Banking and Currency}
The congressional committee whose investigative hearings, led by chief counsel
Ferdinand Pecora and shaped legislatively by Senator Carter Glass, became the public
forum in which Wall Street's boom-era practices were examined and which produced the
era's landmark banking reforms \autocite{libraryjournal1929,dailyeconomy1929}.

\section{Chapter-by-Chapter Narrative Outline}

\emph{As noted above, Sorkin structures the book as a long sequence of date-headed
chapters within two parts rather than as titled, numbered chapters, and a complete
public listing of those individual date-chapters was not found in the sources
consulted \autocite{wikipedia1929book}. The ``chapters'' below are instead a
conceptual reconstruction of the book's narrative movements, grouped in the
chronological order reviewers describe, to help a listener anticipate where the
story is heading without revealing how any individual character's story resolves.}

\begin{enumerate}[leftmargin=1.6em, itemsep=0.6em]

\item \textbf{The Boom in Full Swing (early 1929).} The book opens amid a
roaring bull market. \charname{Charles Mitchell} and \charname{National City Bank}
keep credit flowing to speculators even as the \charname{Federal Reserve} grows
uneasy, setting up the central tension between Wall Street's appetite and
Washington's caution \autocite{airmail1929,supersummary1929}.

\item \textbf{A New President Takes Office (March 1929).} \charname{Herbert Hoover}
is inaugurated into a market he privately mistrusts, while \charname{Andrew Mellon}
continues to shape a Treasury posture of minimal interference in the boom
\autocite{supersummary1929,csmonitor1929}.

\item \textbf{The Machinery of Speculation (spring--summer 1929).} Pool operators
such as \charname{William C.\ Durant}, \charname{John J.\ Raskob}, and
\charname{Michael J.\ Meehan} drive up favored stocks on the \charname{New York
Stock Exchange}, while \charname{Thomas Lamont} and the House of \charname{J.P.\
Morgan \& Co.} watch the market's temperature from above \autocite{supersummary1929}.

\item \textbf{Cracks Beneath the Surface (late summer 1929).} As the market reaches
its historic peak, independent skeptics like \charname{Jesse Livermore} begin to
sense the boom has outrun reality, even as public enthusiasm remains near its height
\autocite{airmail1929,mises1929}.

\item \textbf{The Tremors Begin (early--mid October 1929).} Confidence starts to
waver on the exchange floor; \charname{Richard Whitney} and the leading bankers
grow increasingly attentive to signs of instability \autocite{supersummary1929}.

\item \textbf{Black Thursday (October 24, 1929).} A wave of panic selling hits the
\charname{New York Stock Exchange}; \charname{Thomas Lamont}, \charname{J.P.\
Morgan Jr.}, and their fellow bankers organize a coordinated show of confidence on
the floor, with \charname{Richard Whitney} acting as their public face
\autocite{supersummary1929,csmonitor1929}.

\item \textbf{A Fragile Reprieve (October 25--28, 1929).} A brief stabilization
follows the bankers' intervention, and the country debates whether the worst has
passed \autocite{supersummary1929}.

\item \textbf{Black Tuesday and Its Immediate Wake (October 29, 1929 onward).} Selling
overwhelms the market again, and the scale of the losses becomes clear across Wall
Street, testing every figure introduced so far, from \charname{Charles Mitchell} to
\charname{Herbert Hoover} \autocite{supersummary1929,csmonitor1929}.

\item \textbf{Washington Tries to Steady the Country (late 1929--1930).}
\charname{Herbert Hoover} convenes business leaders and urges confidence, while the
\charname{Federal Reserve System} and figures like \charname{George Harrison}
weigh how aggressively to respond \autocite{supersummary1929}.

\item \textbf{The Economy Contracts (1930--1931).} As economic conditions worsen and
bank failures spread through the financial system, the limits of voluntary,
business-led solutions become increasingly apparent \autocite{supersummary1929}.

\item \textbf{A Nation in Crisis, A New Administration (1932--1933).}
\charname{Franklin D.\ Roosevelt} rises to national prominence and is elected
president amid a deepening banking crisis, setting the stage for a fundamentally
different relationship between Washington and Wall Street \autocite{supersummary1929,
libraryjournal1929}.

\item \textbf{The Reckoning (1933--1934).} \charname{Ferdinand Pecora} leads the
\charname{U.S.\ Senate Committee on Banking and Currency}'s hearings, calling
\charname{Charles Mitchell}, \charname{Richard Whitney}, and other figures from
earlier in the book to account publicly for their conduct during the boom
\autocite{libraryjournal1929,dailyeconomy1929}.

\item \textbf{Rebuilding the Rules (1933--1934).} \charname{Carter Glass} and his
congressional allies, working with the \charname{Roosevelt} administration,
translate the hearings' revelations into the era's landmark banking and securities
reforms, reshaping the relationship between \charname{National City Bank},
\charname{J.P.\ Morgan \& Co.}, and the institutions that regulate them
\autocite{libraryjournal1929}.

\item \textbf{Looking Back.} The book closes by drawing a line from the 1929 crash
and its reforms toward the financial crises of later decades, reflecting on what the
era's central figures and institutions ultimately teach about market manias and
their aftermath \autocite{axios1929,dailyeconomy1929}.

\end{enumerate}

\printbibliography[title={References}]

\end{document}
